\subsection{Continuation approximation and the contribution of NBO}
A reusable continuation is not a new principle of dynamic programming.
Postdecision-state approximation separates a controlled transition from the
subsequent uncertainty and is a standard device in approximate dynamic
programming \citep{powell2011}. Regression can estimate a continuation relevant
to an economic stopping decision, as in \citet{longstaff2001}. Batch fitted
value methods turn transition data into successive supervised regression
problems \citep{ernst2005}; approximate policy iteration likewise separates
an evaluation error from an imperfect improvement step
\citep{scherrer2015}. These methods establish the appropriate comparison
class for NBO. The relevant question is not whether continuation fitting is
possible, but which fitted object, decision rule, and verification account
are useful for the stated economic exercise.

The actor--critic distinction is also established. In
\citet{konda1999}, a critic supplies information for a parameterized policy
update under explicit stochastic-approximation conditions. NBO uses a scalar
policy continuation and its derivative channels in a feasible economic
improvement map. The scalar representation couples value and gradient
predictions through one potential; a separately fitted vector field need
not have this property. That structural distinction is testable but does
not imply that a scalar neural network must outperform a costate surrogate,
a direct actor, or a conventional scalar regression. The menu's vector
baseline and the new quadratic, radial-basis, and shared-actor baselines
therefore remain in the primary comparisons.

The reuse comparison is closely related to simulation optimization. A cached
sample-average objective can be optimized at many current tasks. Reusing
innovations is legitimate; reusing their payoff at a different postdecision
state without reevaluation is not. The new experiment gives cached simulation
and regression the same finite-law information. The regressions share their
value and gradient labels, and the shared actor receives the same task
variables. Full-support enumeration is a particularly strong comparator
when the law has only sixteen atoms. Its tractability is disclosed rather
than suppressed to manufacture an advantage for approximation.

Several mathematical components are standard. Centering removes constants
from action comparisons; the squared-error difference cancels common noise;
shortest-path closure solves difference constraints; and a strong-concavity
bound converts a first-order residual into an objective gap. For stochastic
stopping, time-uniform confidence sequences provide a general route to valid
repeated assessment \citep{howard2021}. We do not claim any of these principles
as a new theorem of dynamic programming, graph theory, or probability. The
contribution developed here is their explicit economic linkage: which
continuation is reusable, how approximation and changes in future primitives
enter decision loss, which computed action is certified, and which work is
actually paid before that certificate is returned. The deterministic finite-law
experiment uses no confidence sequence and spends no additional probability.
Its scalar-interval guarantee and the earlier statistical menu guarantee
are different applications of that linkage.
